% Unit 9 exam -- 25 MCQ + 1 long + 1 short FRQ = 39 pts, 69 min.
\documentclass{apchem-exam}

\apcunit{9}
\apcunittitle{Thermodynamics and Electrochemistry}
\apcdoctitle{Unit Exam}
\apcsubtitle{CED 9.1--9.11 \textbullet\ Fri Mar 19 \textbullet\ 69 minutes \textbullet\ 39 points}

\begin{document}
\apctitleblock

\begin{apcnote}[Format and scope]
Section I: 25 multiple choice, \textbf{no calculator}, 38 minutes.
Section II: one long and one short free response, calculator permitted,
31 minutes.

Per the CED, topic 9.10 is assessed \textbf{without the Nernst equation} ---
predict the direction of a change and justify it with $Q$ against $K$.
Topic 9.8 excludes labelling electrodes as positive or negative; identify
them as anode and cathode.

Write \textbf{thermodynamically favored}, not ``spontaneous''.
\end{apcnote}

\examsection{I}{Multiple Choice}{25 questions \textbullet\ 38 minutes \textbullet\ 1 point each}{\nocalculator}

\begin{mcq}
  Entropy is best described as a measure of
  \choices
    \correct{the dispersal of energy and matter among available
      arrangements}
    \choice{the disorder of the objects in a system}
    \choice{the heat content of a system}
    \choice{the rate at which a reaction proceeds}
\end{mcq}
\why{Entropy counts microscopic arrangements; it is unrelated to rate.}
\distractor{B}{the usual shorthand, but it describes appearance rather than
the quantity measured}

\begin{mcq}
  Which change has the most positive $\Delta S^\circ$?
  \choices
    \choice{\ce{H2O(g) -> H2O(l)}}
    \correct{\ce{CO2(s) -> CO2(g)}}
    \choice{\ce{2NO2(g) -> N2O4(g)}}
    \choice{\ce{Ag+(aq) + Cl^-(aq) -> AgCl(s)}}
\end{mcq}
\why{A solid becoming a gas is the largest possible increase in dispersal.}

\begin{mcq}
  For \ce{N2(g) + 3H2(g) -> 2NH3(g)}, $\Delta S^\circ$ is negative
  principally because
  \choices
    \correct{four moles of gas become two}
    \choice{the reaction is exothermic}
    \choice{ammonia is a polar molecule}
    \choice{the reaction requires a catalyst}
\end{mcq}
\why{When the number of gas moles changes, that term dominates.}

\begin{mcq}
  Which statement about $S^\circ$ is correct?
  \choices
    \correct{$S^\circ$ for \ce{O2(g)} is a positive, nonzero value}
    \choice{$S^\circ$ for any element in its standard state is zero}
    \choice{$S^\circ$ is negative for all stable compounds}
    \choice{$S^\circ$ cannot be measured, only changes in it}
\end{mcq}
\why{The third law fixes a true zero, so absolute entropies exist and
elements have no reason to sit at it.}
\distractor{B}{true of $\Delta H_f^\circ$ and $\Delta G_f^\circ$, not of
$S^\circ$}

\begin{mcq}
  A process has $\Delta S_{\text{sys}} < 0$ and is thermodynamically
  favored. It follows that
  \choices
    \correct{$\Delta S_{\text{surr}}$ is positive and larger in magnitude}
    \choice{the second law has been violated}
    \choice{$\Delta H$ must be positive}
    \choice{$\Delta S_{\text{univ}}$ is negative}
\end{mcq}
\why{The second law constrains the universe, not the system.}

\begin{mcq}
  $\Delta G^\circ = \Delta H^\circ - T\Delta S^\circ$. A reaction with
  $\Delta H^\circ < 0$ and $\Delta S^\circ > 0$ is favored
  \choices
    \correct{at all temperatures}
    \choice{only at high temperature}
    \choice{only at low temperature}
    \choice{at no temperature}
\end{mcq}
\why{Both terms make $\Delta G^\circ$ negative, so $T$ cannot change the
outcome.}

\begin{mcq}
  A reaction with $\Delta H^\circ > 0$ and $\Delta S^\circ > 0$ becomes
  favored
  \choices
    \correct{above $T = \Delta H^\circ/\Delta S^\circ$}
    \choice{below $T = \Delta H^\circ/\Delta S^\circ$}
    \choice{at all temperatures}
    \choice{only if a catalyst is added}
\end{mcq}
\why{The entropy term grows with $T$ until it overtakes the enthalpy term.}
\distractor{D}{a catalyst never changes favorability}

\begin{mcq}
  A student computes $\Delta G^\circ = 178.5 - 298(161)$ and obtains
  $-47{,}800$~kJ. The error is that
  \choices
    \correct{$\Delta S^\circ$ was not converted from J/K to kJ/K}
    \choice{the temperature should be in \si{\celsius}}
    \choice{the sign of $\Delta H^\circ$ is wrong}
    \choice{$\Delta S^\circ$ should have been divided by $T$}
\end{mcq}
\why{The magnitude is the giveaway: no ordinary reaction has a free energy
change of tens of thousands of kJ.}

\begin{mcq}
  For a reaction at equilibrium,
  \choices
    \correct{$\Delta G = 0$}
    \choice{$\Delta G^\circ = 0$}
    \choice{$K = 0$}
    \choice{$\Delta S_{\text{univ}} > 0$}
\end{mcq}
\why{$\Delta G^\circ$ is zero only in the special case $K = 1$.}
\distractor{B}{a very common slip --- $\Delta G^\circ$ is fixed at a given
temperature}

\begin{mcq}
  If $\Delta G^\circ$ for a reaction is positive, then
  \choices
    \correct{$K < 1$}
    \choice{$K > 1$}
    \choice{$K = 0$}
    \choice{the reaction cannot occur under any conditions}
\end{mcq}
\why{$\Delta G^\circ = -RT\ln K$.}
\distractor{D}{a small enough $Q$ still makes $\Delta G$ negative}

\begin{mcq}
  Reaction A has $\Delta G^\circ = -5$~kJ and reaction B has
  $\Delta G^\circ = -50$~kJ. Compared with $K_A$, $K_B$ is
  \choices
    \correct{very much larger}
    \choice{ten times larger}
    \choice{ten times smaller}
    \choice{the same}
\end{mcq}
\why{$K$ depends exponentially on $\Delta G^\circ$.}
\distractor{B}{treats the relationship as linear}

\begin{mcq}
  A reaction has $\Delta G^\circ = -180$~kJ but no measurable rate at
  \SI{298}{\kelvin}. The best explanation is
  \choices
    \correct{a large activation energy}
    \choice{an error in the value of $\Delta G^\circ$}
    \choice{that $K$ must be less than 1}
    \choice{that the reaction is endothermic}
\end{mcq}
\why{Favorability and rate are independent; this is kinetic control.}

\begin{mcq}
  Adding a catalyst to a reaction changes
  \choices
    \correct{the activation energy only}
    \choice{$\Delta G^\circ$ and $K$}
    \choice{$\Delta H^\circ$ only}
    \choice{$K$ but not the rate}
\end{mcq}
\why{A catalyst lowers the barrier without moving the reactant or product
free energies.}

\begin{mcq}
  \ce{NH4NO3} dissolves endothermically yet is highly soluble. This is
  because
  \choices
    \correct{$T\Delta S^\circ$ exceeds $\Delta H^\circ$}
    \choice{the lattice energy is negligible}
    \choice{$\Delta S^\circ$ is negative}
    \choice{dissolving is always favored}
\end{mcq}
\why{Entropy drives this dissolution against the enthalpy term.}

\begin{mcq}
  When an ionic solid dissolves, the reorganization of solvent molecules
  into hydration shells contributes
  \choices
    \correct{a negative $\Delta S$}
    \choice{a positive $\Delta S$}
    \choice{a positive $\Delta H$}
    \choice{nothing to either term}
\end{mcq}
\why{Ordering the solvent around ions reduces the number of arrangements
available to it.}

\begin{mcq}
  An unfavorable reaction is coupled to a favorable one. For the coupled
  process to be favored, the two reactions must share a common intermediate
  and
  \choices
    \correct{the sum of the two $\Delta G^\circ$ values must be negative}
    \choice{both reactions must be exothermic}
    \choice{the favorable reaction must be faster}
    \choice{the two $\Delta G^\circ$ values must be equal in magnitude}
\end{mcq}
\why{Free energies add.}

\begin{mcq}
  In any electrochemical cell, oxidation occurs at the
  \choices
    \correct{anode}
    \choice{cathode}
    \choice{salt bridge}
    \choice{negative electrode in every case}
\end{mcq}
\why{An~Ox holds in galvanic and electrolytic cells alike.}
\distractor{D}{the sign reverses between cell types, which is why the CED
excludes it}

\begin{mcq}
  The purpose of a salt bridge is to
  \choices
    \correct{carry ions so both solutions stay electrically neutral}
    \choice{carry electrons between the half-cells}
    \choice{prevent the two solutions from reacting}
    \choice{increase the cell potential}
\end{mcq}
\why{Without it, charge builds up and electron flow stops almost at once.}
\distractor{B}{electrons travel through the wire, not the bridge}

\begin{mcq}
  An electrolytic cell has
  \choices
    \correct{$E^\circ_{\text{cell}} < 0$ and $\Delta G^\circ > 0$}
    \choice{$E^\circ_{\text{cell}} > 0$ and $\Delta G^\circ < 0$}
    \choice{$E^\circ_{\text{cell}} > 0$ and $\Delta G^\circ > 0$}
    \choice{$E^\circ_{\text{cell}} = 0$ and $\Delta G^\circ = 0$}
\end{mcq}
\why{It requires an external supply, which is CED 9.7's ``external energy''
in hardware.}
\distractor{D}{that describes a dead battery}

\begin{mcq}
  Given \ce{Zn^2+}/\ce{Zn} $= -0.76$~V and \ce{Cu^2+}/\ce{Cu} $= +0.34$~V,
  $E^\circ_{\text{cell}}$ for \ce{Zn + Cu^2+ -> Zn^2+ + Cu} is
  \choices
    \correct{$+1.10$~V}
    \choice{$-1.10$~V}
    \choice{$-0.42$~V}
    \choice{$+0.42$~V}
\end{mcq}
\why{$E^\circ = 0.34 - (-0.76)$.}
\distractor{D}{adds the two potentials instead of subtracting}

\begin{mcq}
  For \ce{Cu + 2Ag+ -> Cu^2+ + 2Ag}, with \ce{Ag+}/\ce{Ag} $= +0.80$~V and
  \ce{Cu^2+}/\ce{Cu} $= +0.34$~V, $E^\circ_{\text{cell}}$ is
  \choices
    \correct{$+0.46$~V}
    \choice{$+1.26$~V}
    \choice{$+1.94$~V}
    \choice{$-0.46$~V}
\end{mcq}
\why{Doubling the silver half-reaction does not double its potential ---
$E^\circ$ is intensive.}
\distractor{B}{doubles $E^\circ$ along with the coefficients, the classic
error}

\begin{mcq}
  Two cells have the same $E^\circ_{\text{cell}}$ but transfer 1 and 3
  electrons respectively. Compared with the first, the second cell has
  \choices
    \correct{the same voltage and three times the magnitude of
      $\Delta G^\circ$}
    \choice{three times the voltage and the same $\Delta G^\circ$}
    \choice{three times both}
    \choice{the same value of both}
\end{mcq}
\why{Potential is intensive; free energy is an amount of energy and scales
with $n$.}

\begin{mcq}
  A galvanic cell has been running for some time. Compared with its initial
  state,
  \choices
    \correct{$Q$ has increased and $E_{\text{cell}}$ has decreased}
    \choice{$Q$ has decreased and $E_{\text{cell}}$ has increased}
    \choice{$K$ has decreased}
    \choice{$E^\circ_{\text{cell}}$ has decreased}
\end{mcq}
\why{The reaction moves toward equilibrium; $K$ and $E^\circ$ are fixed at
a given temperature.}
\distractor{D}{$E^\circ$ is a standard-state constant, unlike
$E_{\text{cell}}$}

\begin{mcq}
  For \ce{Zn(s) + Cu^2+(aq) -> Zn^2+(aq) + Cu(s)}, increasing
  $[\ce{Cu^2+}]$ will
  \choices
    \correct{increase $E_{\text{cell}}$, because $Q$ decreases}
    \choice{decrease $E_{\text{cell}}$, because $Q$ decreases}
    \choice{increase $E_{\text{cell}}$, because $E^\circ$ increases}
    \choice{leave $E_{\text{cell}}$ unchanged}
\end{mcq}
\why{A reactant concentration appears in the denominator of $Q$; a smaller
$Q$ means a larger displacement from equilibrium.}

\begin{mcq}
  Passing one mole of electrons through solutions of \ce{Ag+},
  \ce{Cu^2+}, and \ce{Al^3+} deposits
  \choices
    \correct{1 mol \ce{Ag}, 0.5 mol \ce{Cu}, and $1/3$ mol \ce{Al}}
    \choice{one mole of each metal}
    \choice{1 mol \ce{Ag}, 2 mol \ce{Cu}, and 3 mol \ce{Al}}
    \choice{amounts that depend on the current used}
\end{mcq}
\why{The half-reaction supplies the electron ratio, read off the ion's
charge.}
\distractor{C}{multiplies where it should divide}

\examsection{II}{Free Response}{2 questions \textbullet\ 31 minutes \textbullet\ calculator permitted}{}

% ---- FRQ 1 (long) ----------------------------------------------------------
\frq{long}{10}{CED 9.8, 9.9, 9.10 \textbullet\ a galvanic cell end to end}

A galvanic cell is assembled from a nickel electrode in \SI{1.0}{\Molar}
\ce{Ni(NO3)2} and a silver electrode in \SI{1.0}{\Molar} \ce{AgNO3}, joined
by a salt bridge.

\medskip
$E^\circ$: \ce{Ag+ + e^- -> Ag}, $+0.80$~V \textbullet\
\ce{Ni^2+ + 2e^- -> Ni}, $-0.23$~V \textbullet\
$F = \SI{96485}{\coulomb\per\mole}$

\begin{parts}
  \item Identify the anode and the cathode, and give the direction of
        electron flow through the wire.
        \answerlines[2]{Nickel is the anode (oxidation --- it has the more
        negative reduction potential); silver is the cathode (reduction).
        Electrons flow through the wire from the nickel electrode to the
        silver electrode.}
  \item Write the balanced overall cell reaction and state $n$.
        \answerlines[2]{\ce{Ni(s) + 2Ag+(aq) -> Ni^2+(aq) + 2Ag(s)};
        $n = 2$.}
  \item Calculate $E^\circ_{\text{cell}}$. \workspace[1.8cm]
        \keyonly{\begin{apcanswer}$E^\circ = 0.80 - (-0.23) =
        \mathbf{+1.03}$~V\end{apcanswer}}
  \item Calculate $\Delta G^\circ$ for the cell reaction, in kJ.
        \workspace[2.4cm]
        \keyonly{\begin{apcanswer}$\Delta G^\circ = -nFE^\circ =
        -(2)(96485)(1.03) = -198{,}759~\text{J} =
        \mathbf{-199}~\text{kJ}$\end{apcanswer}}
  \item The \ce{AgNO3} solution is diluted to \SI{0.010}{\Molar} while the
        nickel half-cell is unchanged. Predict the effect on
        $E_{\text{cell}}$ and justify with $Q$. Do not calculate.
        \answerlines[3]{\ce{Ag+} is a reactant and appears in the
        denominator of $Q$, so lowering its concentration increases $Q$.
        A larger $Q$ places the system nearer equilibrium, reducing the
        driving force, so $E_{\text{cell}}$ decreases.}
  \item State the mass change at each electrode as the cell runs, and give
        $E_{\text{cell}}$ when the cell is exhausted.
        \answerlines[2]{The nickel anode loses mass as it is oxidized; the
        silver cathode gains mass as \ce{Ag+} plates onto it. When
        exhausted, $Q = K$ and $E_{\text{cell}} = 0$.}
\end{parts}

\begin{rubric}
  \rpoint{2}{(a) 1 pt both electrodes correctly identified with oxidation
    at nickel; 1 pt electron flow stated anode $\to$ cathode. Labelling an
    electrode $+$ or $-$ is neither required nor penalized}
  \rpoint{2}{(b) 1 pt balanced equation with \ce{Ag+} doubled; 1 pt
    $n = 2$ --- an unbalanced equation earns 0 for both}
  \rpoint{1}{(c) $+1.03$~V --- using $2(0.80)$ earns 0}
  \rpoint{2}{(d) 1 pt substitution into $-nFE^\circ$ with $n = 2$; 1 pt
    $-199$~kJ with the conversion from joules --- reporting
    $-198{,}759$~kJ earns 0 for this point}
  \rpoint{2}{(e) 1 pt states $Q$ increases because \ce{Ag+} is a reactant;
    1 pt concludes $E_{\text{cell}}$ decreases. A bare ``decreases''
    with no $Q$ reasoning earns 0}
  \rpoint{1}{(f) both mass changes AND $E_{\text{cell}} = 0$ required}
\end{rubric}

% ---- FRQ 2 (short) ---------------------------------------------------------
\frq{short}{4}{CED 9.3, 9.4 \textbullet\ favorability and rate}

For \ce{CaCO3(s) -> CaO(s) + CO2(g)}, $\Delta H^\circ = +178.5$~kJ and
$\Delta S^\circ = +161$~\si{\joule\per\kelvin}.

\begin{parts}
  \item Calculate $\Delta G^\circ$ at \SI{298}{\kelvin} and state whether
        the reaction is thermodynamically favored. \workspace[2.4cm]
        \keyonly{\begin{apcanswer}$\Delta G^\circ = 178.5 - 298(0.161)
        = 178.5 - 48.0 = \mathbf{+131}$~kJ. Positive, so \textbf{not}
        thermodynamically favored at \SI{298}{\kelvin}.\end{apcanswer}}
  \item Calculate the temperature above which the reaction becomes favored.
        \workspace[1.8cm]
        \keyonly{\begin{apcanswer}$T = \Delta H^\circ/\Delta S^\circ =
        178.5/0.161 = \mathbf{1109}~\si{\kelvin}$\end{apcanswer}}
  \item A student claims that adding a catalyst would allow the
        decomposition to proceed at \SI{298}{\kelvin}. Evaluate the claim.
        \answerlines[3]{The claim is wrong. A catalyst lowers the
        activation energy and so increases the rate, but it does not change
        $\Delta H^\circ$, $\Delta S^\circ$, $\Delta G^\circ$, or $K$. Since
        $\Delta G^\circ$ is positive at \SI{298}{\kelvin}, the reaction is
        not favored and no catalyst can make it so.}
\end{parts}

\begin{rubric}
  \rpoint{2}{(a) 1 pt correct substitution with $\Delta S^\circ$ converted
    to kJ/K; 1 pt $+131$~kJ and ``not favored''. An answer near
    $-47{,}800$ signals the unconverted entropy and earns 0 for both}
  \rpoint{1}{(b) $1109$~K (accept \SI{1100}{\kelvin} to 2 s.f.)}
  \rpoint{1}{(c) rejects the claim AND states that a catalyst leaves
    $\Delta G^\circ$ or $K$ unchanged --- ``no'' alone earns 0}
\end{rubric}

\examtotal

\end{document}
